\subsection{RGA 二维加速}\label{sec:rga}

机器人采集到的每一帧并不能直接送入检测网络。相机交付的是 YUV 排布的像素，而 YOLOv8 的输入要求是连续的 RGB 平面，尺寸还需缩放到模型的输入分辨率。若交由 CPU 逐像素完成色彩空间转换与缩放，一帧的开销与一次推理相当，感知与计算彼此争抢同一个核，流水线随即被这一步拖住。RK3588 为此集成了 RGA 二维栅格图形加速器，它在硬件中完成色彩空间转换、缩放与图层拼接，并能直接读写物理内存，使这一步从计算路径上让出。

\repfig{rga-architecture.png}{RGA 二维加速的软件栈与数据通路，librga 经 \texttt{/dev/rga} 提交 \texttt{rga\_req}，内核翻译为命令块交由 RGA2 引擎执行，结果直接写入下游缓冲}{fig:rga}

RGA 的软件入口在 Linux 上是 \texttt{/dev/rga} 与用户态的 librga。应用以一份 \texttt{rga\_req} 描述源与目标的地址、格式、尺寸与裁剪窗口，经一次 \texttt{RGA\_BLIT\_SYNC} 同步阻塞调用提交，内核把它翻译为一段命令块，交由 RGA2 引擎执行。本队在 Starry 上重建了这条通路，一个平台无关的 \texttt{\#![no\_std]} 驱动核心负责命令块编码与寄存器时序，\texttt{/dev/rga} 设备节点对齐 librga 的 ioctl 约定，使既有的 librga 二进制无需改动即可驱动这块硬件。

\paragraph{对齐用户态的二进制接口}
librga 与内核之间并非函数调用，而是一份按字节对齐的结构体经 ioctl 传递，任何一处字段错位都会让引擎读到错误的地址或尺寸。\texttt{rga\_req} 在 LP64 下恰为 504 字节，其中嵌入三份各 56 字节的图像描述符，缓冲导入用的 \texttt{rga\_external\_buffer} 为 288 字节，这些尺寸以编译期断言固定下来，任何字段增删都会在编译时暴露。早期的镜像漏掉了 \texttt{rotate\_mode}、\texttt{rd\_mode} 等几个较新的字段，使其后所有字段整体偏移，引擎据此取到错位的旋转矩阵与地址。librga 在打开设备时还会查询驱动版本与硬件版本以选择匹配的结构体布局，因此设备节点须如实报告 MultiRGA v1.3.1 与 RK3588 上 RGA2 核的版本号，librga 才会按这一布局组包并把请求投向 RGA2 而非另一代核心。

\paragraph{在三个核中找到正确的那一个}
RK3588 的设备树把 RGA 暴露为三个独立设备，两个 RGA3 核与一个 RGA2 核。本流水线所需的色彩转换与缩放走 RGA2 路径，若按设备枚举顺序取首个核，落到的是一个没有 RGA2 引擎的 RGA3 核，提交随即以无此设备失败。设备节点因此按版本而非按枚举序检索，锁定真正承载 RGA2 的那个核再提交。

\paragraph{让引擎真正跑起来并报告完成}
RGA2 的执行时序对寄存器写入的顺序敏感。引擎须先被置入主模式，命令起始信号才会触发一次命令 DMA 取指，停留在从模式时起始信号什么也取不到。完成信号同样有讲究，表示全部命令完成的状态位只有在对应的中断使能于起始之前被置位时才会锁存。最初省略了这一步，引擎明明已经跑完，完成位却始终为零，轮询一路超时，看上去像是硬件挂死。据厂商时序在主模式与起始之间补上这次对中断使能位的读改写后，完成位正常锁存，既有的忙等轮询即可观察到它，无需真正布线中断。命令缓冲的基址按原始字节地址写入，不做页表基址那样的右移，这一区别在 MMU 关闭、直接使用物理地址的本路径上是必须的。

\paragraph{一个把缩放挂死的定点系数}
缩放系数以 16.16 定点表示。缩小时正确的系数是目标维度比源维度，其值不超过一，写入低半字；最初误用了它的倒数，缩小时系数大于一，使源指针每步跨过有效区域，扫描器随即停在忙状态约五十毫秒不再前进。把缩小与放大两种情形分别按厂商参数计算后，缩放回到一次提交即完成。

\paragraph{让三个加速器读同一段内存}
RGA 在本流水线里不孤立工作，它读上一级解码的输出，又把结果直接写入 NPU 的输入缓冲，这一共享缓冲的机制见第~\ref{sec:dmaheap}~节。RGA2 在本阶段以 MMU 关闭、物理直通的方式运行，\texttt{/dev/rga} 收到 dma-buf 文件描述符后解析出物理地址交给引擎，并在提交与轮询期间持有该分配的引用计数，使另一线程的并发释放不会在操作中途抽走内存，提交前后还须显式完成缓存与设备方向的同步，引擎与 CPU 才看到一致的数据。

\textbf{驱动的五项基本操作在开发板上以逐像素 CRC 校验通过自测，该工作以合并请求 \href{https://github.com/rcore-os/tgoskits/pull/1388}{\#1388} 提交至上游。}在 Linux 上的受控对照中，以 RGA 替代 CPU 缩放可将该阶段由 15.82\,ms 降至 1.72\,ms，约九倍；Starry 接入这条路径后，色彩空间转换的 p50 为 0.776\,ms，缩放也接近归零。这一步把图像准备从计算路径上彻底让出，为后续把推理放置到大核腾出了余量，相关数据见第~\ref{sec:opt}~节。
